\paragraph{A current-STA window exposes genuine endpoint-only progress.}
The previous exact-worst selector has no improving single-connection candidate at the common R1-derived 5.854659868 ns state. We expand the current slack window to 0.2 ns and temporarily restrict the operator to single-sink repairs, isolating selection from joint-scope changes. Both acceptance policies start at the same previously automatic accepted state; this is a continuation experiment, not a new cold-start benchmark. The first window selects 51 connections. Existing exact queries are reused only with matching retained trees, anchors, occupancy, graph and query implementation.

Strict CPD rejects four legal single-connection candidates. The endpoint-vector branch accepts repairs to 7647/pin~19, 7647/pin~32, 5487/pin~24 and 5487/pin~26 in sequence, re-running target selection on each accepted state's native STA. Each transition is a strict-CPD rejection under the exact same before/after numbers, but improves respectively one, one, two and one setup components without regression. In total five distinct endpoint/domain components improve; hold-violation guards pass. The fixed initial 0.2 ns endpoint population decreases from fourteen to eleven. CPD remains exactly 5.854659868020917 ns, and the worst connection group remains 76/pin~1 and 48/pin~1. Thus these steps reduce near-critical endpoint pressure but have not yet lowered CPD or removed the current worst bottleneck.

The endpoint branch uses 44 SB pairs, 972 CB sources and planar wirelength 218247, versus 41, 971 and 218171 initially. The timing-vector gain therefore costs three additional SB pairs, one CB source and 76 wirelength units. The complete shared-cache two-policy experiment takes 595.235 s: 314.620 s for candidate construction/query/audit wrappers and 271.105 s for candidate evaluation wrappers, with the remainder in orchestration and decision bookkeeping. It includes 238 new exact capacity queries, three reused queries and seven new FULL/STA evaluations; the first accepted endpoint step reuses the strict branch's identical evaluation. These are development-flow costs, not independent cold-start branch runtimes. Exhausting this single-sink window is a local result only. Whether the accepted plateau states enable a subsequent joint repair must be tested with the same successor opportunities from both the original strict state and the endpoint branch state.

\paragraph{Matched native opportunity controls separate attribution.}
Native VPR connection searches start from the exact state before each accepted endpoint step, with the same single-sink release, retained branches, full-device domain and acceptance opportunities. Three steps generate exactly the same complete physical trees as the capacity-query candidates, including all parent edges, switches and sink identities. Their already certified FULL/STA results are reused only after this identity proof and matching packing, placement, graph and applied constraints. The native 5487/pin~24 step instead leaves all physical connection delays unchanged; a separate fresh FULL/STA confirms no endpoint-vector gain, so both policies reject it. The capacity candidate's two improvements are approximately 7.638 ps each. The current native cost mixes criticality-weighted delay and congestion, whereas the capacity oracle minimizes physical delay subject to feasibility. This remaining discrepancy therefore does not establish a superior search engine; a same-objective native feasibility-filter control is required. The four actual native searches use eleven pressure attempts in total, with core transaction times from 0.020 to 0.216 s. Incremental control orchestration takes 146.584 s, including exact-equivalence checks, the one differing candidate's fresh evaluation and three explicit prior-evaluation reuses. This is a matched-state attribution test, not an autonomous native branch after divergence or a formal speed comparison.
